% !TEX root = 09.tex
\documentclass[a4paper,10pt]{article}

\usepackage{pdfsync}
\usepackage[utf8]{inputenc}
\usepackage{microtype}
\usepackage{mathtools}
\usepackage{amsthm,amsfonts,amsmath,amssymb}
\usepackage{mathabx}
% \usepackage{MnSymbol}
\usepackage{hyperref}
\usepackage{bbold}
\usepackage[textwidth=2cm,textsize=small]{todonotes}
\usepackage{subcaption}
\usepackage{enumerate}
\usepackage{xspace}
\usepackage{stmaryrd}
\usepackage{verbatim}
\usepackage{cleveref}

\hypersetup{
    colorlinks,
    linkcolor={red!50!black},
    citecolor={blue!50!black},
    urlcolor={blue!80!black}
}

\author{}

\newcommand{\ignore}[1]{}

\newcommand{\st}{s.t.}
\newcommand{\ie}{i.e.}
\newcommand{\eg}{e.g.}

\newcommand{\sep}{\;|\;}
\renewcommand{\rule}[2]{\dfrac{#1}{#2}}
\newcommand{\set}[1]{\left\{#1\right\}}
\newcommand{\setof}[2]{\set{#1 \;\middle|\; #2}}
\newcommand{\tuple}[1]{\left( #1 \right)}
\newcommand{\pair}[2]{\tuple {#1, #2}}
\newcommand{\closure}[3]{\tuple {#1, #2, #3}}
\newcommand{\sem}[1]{\llbracket{#1}\rrbracket}
\newcommand{\powerset}[1]{\mathcal P(#1)}
\newcommand{\map}[2]{\mathsf{map}\; #1\; #2}
\newcommand{\filter}[2]{\mathsf{filter}\; #1\; #2}
\newcommand{\floor}[1]{\left\lfloor #1 \right\rfloor}
\newcommand{\PreRel}[2]{\mathsf{Pre}_{#1}(#2)}
\newcommand{\Pre}[1]{\PreRel {} {#1}}
\newcommand{\PreStar}[1]{\mathsf{Pre}^*(#1)}
\newcommand{\Post}[1]{\mathsf{Post}(#1)}
\newcommand{\PostStar}[1]{\mathsf{Post^*}(#1)}
\newcommand{\Conf}{\mathsf{Conf}}
\newcommand{\upwardclosure}[1]{#1\uparrow}
\newcommand{\card}[1]{\left| #1 \right|}
\newcommand{\lang}[1]{L(#1)}
\newcommand{\goesto}[1]{\xrightarrow{#1}}
\newcommand{\e}{\varepsilon}

\newcommand{\N}{\mathbb N}
\newcommand{\Z}{\mathbb Z}
\newcommand{\Q}{\mathbb Q}

% LTL
\newcommand{\true}{\mathbf{true}}
\newcommand{\false}{\mathbf{false}}
\newcommand{\X}{\mathbf X}
\newcommand{\U}{\mathbf U}
\newcommand{\Release}{\mathbf R}
\newcommand{\F}{\mathbf F}
\newcommand{\G}{\mathbf G}

% CTL
\newcommand{\EX}{\mathbf{EX}}
\newcommand{\AX}{\mathbf{AX}}
\newcommand{\EF}{\mathbf{EF}}
\newcommand{\AF}{\mathbf{AF}}
\newcommand{\EG}{\mathbf{EG}}
\newcommand{\AG}{\mathbf{AG}}
\newcommand{\EU}[2]{\mathbf E(#1 \U #2)}
\newcommand{\AU}[2]{\mathbf A(#1 \U #2)}
\newcommand{\ER}[2]{\mathbf E(#1 \Release #2)}
\newcommand{\AR}[2]{\mathbf A(#1 \Release #2)}

\newcommand{\NL}{\textsf{NL}}
\newcommand{\PTIME}{\textsf{PTIME}}
\newcommand{\PSPACE}{\textsf{PSPACE}}

\makeatletter

\def\dasharrowfill@#1#2#3#4{%
        $\m@th
        \thickmuskip0mu
        \medmuskip\thickmuskip
        \thinmuskip\thickmuskip
        \relax
        #4#1\mkern2mu
        \xleaders\hbox{$#4\mkern2mu#2\mkern2mu$}\hfill
        \mkern2mu
        #3$%
}

\def\dashleftarrowfill@{\dasharrowfill@\leftarrow\relbar\relbar}
\def\dashrightarrowfill@{\dasharrowfill@\relbar\relbar\rightarrow}
\def\dashleftrightarrowfill@{\dasharrowfill@\leftarrow\relbar\rightarrow}
\def\dashLeftarrowfill@{\dasharrowfill@\Leftarrow\Relbar\Relbar}
\def\dashRightarrowfill@{\dasharrowfill@\Relbar\Relbar\Rightarrow}
\def\dashLeftrightarrowfill@{\dasharrowfill@\Leftarrow\Relbar\Rightarrow}

\providecommand*\xdashleftarrow[2][]{%
  \ext@arrow 0055{\dashleftarrowfill@}{#1}{#2}}
\providecommand*\xdashrightarrow[2][]{%
  \ext@arrow 0055{\dashrightarrowfill@}{#1}{#2}}
\providecommand*\xdashleftrightarrow[2][]{%
  \ext@arrow 0055{\dashleftrightarrowfill@}{#1}{#2}}
\providecommand*\xdashLeftarrow[2][]{%
  \ext@arrow 0055{\dashLeftarrowfill@}{#1}{#2}}
\providecommand*\xdashRightarrow[2][]{%
  \ext@arrow 0055{\dashRightarrowfill@}{#1}{#2}}
\providecommand*\xdashLeftrightarrow[2][]{%
  \ext@arrow 0055{\dashLeftrightarrowfill@}{#1}{#2}}

\def\slashedarrowfill@#1#2#3#4#5{%
$\m@th\thickmuskip0mu\medmuskip\thickmuskip\thinmuskip\thickmuskip
  \relax#5#1\mkern-7mu%
  \cleaders\hbox{$#5\mkern-2mu#2\mkern-2mu$}\hfill
  \mathclap{#3}\mathclap{#2}%
  \cleaders\hbox{$#5\mkern-2mu#2\mkern-2mu$}\hfill
  \mkern-7mu#4$%
}

\def\rightslashedarrowfill@{%
  \slashedarrowfill@\relbar\relbar{\raisebox{1.2pt}{$\scriptscriptstyle\diagup$}}\rightarrow}
\newcommand\xslashedrightarrow[2][]{%
  \ext@arrow 0055{\rightslashedarrowfill@}{#1}{#2}}
\makeatother

\newtheorem{exercise}{Exercise}

\newenvironment{solution}{\begin{proof}[\protect{\textnormal{\emph{Solution.}}}]}{\end{proof}}

\let\oldqedhere\qedhere

\newif\ifstartedinmathmode
\renewcommand*{\qedhere}{%
  \relax\ifmmode\startedinmathmodetrue\else\startedinmathmodefalse\fi%
  \ifstartedinmathmode\tag*{\protect\oldqedhere}\else\ifinlist{\hfill\oldqedhere}{\oldqedhere}\fi%
}

\crefname{equation}{}{}
\crefname{exercise}{Exercise}{Exercises}


\title{Formal verification - Tutorial 09}
\date{18/05/2026}

\begin{document}

\maketitle

\section*{Gale-Stewart games}

\begin{exercise}
    Consider Gale-Stewart games with context-free winning conditions.
    \begin{itemize}
        % \item Are they determined?
        \item Can we decide whether a given player has a winning strategy?
    \end{itemize}
\end{exercise}
\begin{solution}
    No, GS games with $\omega$-context-free winning conditions are undecidale.
    We can reduce from the universality problem for $\omega$-context-free languages, which is known to be undecidable.
    For instance, consider an $\omega$-CFL $L \subseteq \Sigma^\omega$.
    Consider the GS game where P1 plays letters from $\Sigma$ and P2 plays letters from a singleton alphabet $\set{a}$.
    Consider the winning condition for the second player
    $W_{II} := \setof {a_0 \$ a_1 \$ a_2 \$ \cdots} {a_0 a_1 a_2 \cdots \in L}$.
    %
    By the closure properties of $\omega$-CFLs, $W_{II}$ is an $\omega$-CFL.
    %
    Then P2 has a winning strategy iff $L = \Sigma^\omega$.

    A similar argument can be used to show undecidability when considering a $\omega$-CFL winning condition $W_I$ for P1.
\end{solution}

% \begin{exercise}
%     Let's allow Player II a fixed lookahead on the moves of Player I.
%     Is it decidable whether Player II has a winning strategy?
% \end{exercise}
% \begin{solution}
% \end{solution}

% \begin{exercise}
%     Let's allow Player II a fixed \emph{negative} lookahead on the moves of Player I.
%     Is it decidable whether Player II has a winning strategy?
% \end{exercise}
% \begin{solution}
% \end{solution}

% \begin{exercise}
%     Consider MSO sublogics such as $FO(\leq)$ and $FO(+1)$
%     and corresponding GS games.
%     %
%     Are winning strategies definable already in these sublogics?
% \end{exercise}
% \begin{solution}
% \end{solution}

\section{Closure under uniformisation}

\begin{exercise}[\protect{\cite[Sec.~6]{Thomas:FOSSACS:2009}}]
    Consider Presburger arithmetic $\tuple{\N, +, \leq}$.
    %
    Construct a formula $\varphi(X, Y)$ with two monadic free variables $X, Y$
    \st~$\exists Y \forall X \varphi(X, Y)$ is true, but
    there is no \emph{recursive} witness $Y = \tilde Y$ \st~$\forall X \varphi(X, \tilde Y)$ is true.
\end{exercise}
\begin{solution}
    Let $\varphi_{\text{sq}}(X)$ say that $X$ is the set of squares $S := \set{0^2, 1^2, 2^2, 3^2, \dots}$.
    %
    It is definable in Presburger arithmetic since the distance of three subsequent squares
    $a = n^2$, $b = (n+1)^2$, $c = (n+2)^2$ increases by two:
    $(c - b) - (b - a) = 2$.
    %
    We can then write $\varphi_{\text{sq}}(X)$ as the conjunction of the following three formulas:
    %
    \begin{itemize}
        \item $X(0) \land X(1)$,
        \item $X$ is infinite, and
        \item ``the distance of three consecutive elements of $X$ increases by two'' (as described above).
    \end{itemize}
    %
    Since multiplication is definable in the structure $\tuple{\N, +, S}$,
    there exits a non-recursive set $M \subseteq \N$
    definable by a formula $\psi(S, y)$ over $\tuple{\N, +, S}$ with one free variable $y$.
    %
    Now consider the following Presburger formula
    %
    \begin{align*}
        \varphi(X, Y) \;\equiv\; \varphi_{\text{sq}}(X) \to \forall y. (Y(y) \leftrightarrow \psi(X, y)).
    \end{align*}
    %
    Since there is a unique set of squares, the only solution to the formula above is $Y = M$,
    which is not recursive.
\end{solution}

\section{Deterministic separability}

Consider language $L, M, S \subseteq \Sigma^\omega$.
We say that $S$ \emph{separates} $L$ and $M$ if $L \subseteq S$ and $M \cap S = \emptyset$.

\begin{exercise}
    Show that the following problem is decidable.
    %
    In input we are given two nondeterministic Büchi automata $A, B$ over the same alphabet $\Sigma$
    and a finite set of priorities $P \subseteq \N$,
    and we need to decide whether there is a deterministic parity automaton $S$ over $\Sigma$ with priorities in $P$
    \st~$L(S)$ separates $L(A)$ and $L(B)$.
\end{exercise}
% \begin{solution}
%     TODO.
% \end{solution}

\bibliographystyle{plainurl}
\bibliography{bibliography}

\end{document}